\problema{01 Matrix}{542}{src/01-grafos/542-01-matrix.cpp}{M}

Nos dan una matriz de \texttt{0} y \texttt{1}. Para cada celda debemos calcular
su distancia (en pasos horizontales o verticales) al \texttt{0} más cercano.
La distancia entre celdas adyacentes es 1.

\paragraph{La idea, con un ejemplo de la vida real.} Para cada casilla queremos
saber a cuántos pasos está del \texttt{0} más cercano, moviéndonos solo en
horizontal o vertical (un paso entre casillas pegadas vale 1).

Lo primero que se te ocurre es: «para cada \texttt{1}, salgo a buscar el cero
más cercano». Pero eso es muy lento, porque repites la búsqueda una y otra vez.
El truco es mirarlo al revés: en lugar de que cada \texttt{1} salga a buscar,
hacemos que \emph{todos} los ceros suelten una onda al mismo tiempo, como
cuando tiras muchas piedritas a un charco a la vez. Cada minuto la onda avanza
un paso. La \emph{primera} onda que toca a un \texttt{1} viene, por fuerza, de
su cero más cercano, y el minuto en que lo toca es justo su distancia. Esta
técnica de soltar ondas por capas desde muchos puntos de partida se llama
\emph{BFS multi-fuente} (\emph{BFS} es explorar por capas; \emph{multi-fuente}
es que arrancamos desde varios lugares a la vez).

\begin{center}
\ttfamily
\begin{tabular}{ccc}
0 & 0 & 0\\
0 & 1 & 0\\
1 & 1 & 1\\
\end{tabular}
\qquad$\longrightarrow$\qquad
\begin{tabular}{ccc}
0 & 0 & 0\\
0 & 1 & 0\\
1 & 2 & 1\\
\end{tabular}
\end{center}

\noindent Los ceros valen 0; el \texttt{1} del centro tiene un cero pegado
(distancia 1); la celda de en medio de la última fila está a 2 pasos.

\paragraph{Cómo lo resuelve el código, paso a paso.} Usamos una \emph{cola}:
una fila de espera donde el primero que entra es el primero que sale.
\begin{enumerate}
  \item Recorremos la matriz. A cada casilla con \texttt{0} le ponemos
        distancia 0 y la metemos a la cola (son los puntos de donde salen las
        ondas). A cada \texttt{1} le ponemos distancia «infinito», que
        significa «todavía no sé, muy lejos».
  \item Procesamos la cola. Sacamos una casilla y miramos sus cuatro vecinas.
        Si a una vecina le podemos dar una distancia mejor (la distancia de la
        casilla actual más 1), se la ponemos y la metemos a la cola para que
        siga propagando la onda.
  \item Como cada casilla mejora su distancia una sola vez y luego ya no se
        vuelve a tocar, el recorrido es rápido y ordenado.
\end{enumerate}

\lstinputlisting{src/01-grafos/542-01-matrix.cpp}

\complejidad{$O(m\cdot n)$}{$O(m\cdot n)$}

\paragraph{¿Qué significa esa complejidad?} Aquí $m$ es el número de filas y
$n$ el de columnas, así que $m\cdot n$ es el total de casillas. Que el tiempo
sea $O(m\cdot n)$ significa que cada casilla entra y sale de la cola una sola
vez: tocamos cada una una vez. Comparado con la idea lenta de buscar el cero
por separado para cada \texttt{1}, esto es muchísimo más rápido.

\paragraph{Consejos para la entrevista (Amazon).} La frase mágica que quieren
oír es \emph{BFS multi-fuente}: darte cuenta de que soltar todas las ondas a la
vez convierte un problema que parecía lentísimo en uno rápido. Como alternativa
puedes mencionar la \emph{programación dinámica} (resolver algo grande a partir
de resultados pequeños ya calculados) con dos pasadas: una de arriba-izquierda
hacia abajo y otra de abajo-derecha hacia arriba. También es rápida y no
necesita cola.
